\question[8]
In einer Datenbank sollen Informationen zu Büchern, ihren Kapiteln und den jeweiligen Autoren
gespeichert werden. Hierzu wird zunächst das folgende relationale Datenbankschema definiert:

Dabei gilt: Ein \enquote{Kapitel} gehört genau zu einem Buch. Ein \enquote{Buch} hat im
Ausgangsmodell genau einen Autor.

\begin{itemize}
\item Kapitel ( \pk{KID}, Überschrift, \fk{BID} )
\item Buch ( \pk{BID}, Titel, \fk{AID}, Erscheinungsjahr )
\item Autor ( \pk{AID}, Name, Geburtsjahr )
\end{itemize}

\begin{enumerate}[label=\alph*)]
\item Erläutere anhand dieser Relationen die Begriffe \enquote{Primärschlüssel} und
\enquote{Fremdschlüssel}.

\item Entwirf ein ER-Diagramm zum angegebenen Datenbankschema. Gib auch
die Kardinalitäten aller vorkommenden Beziehungen an.

\item Jemand schlägt vor, das Datenbankschema zu ändern und nur eine
einzige Tabelle zu verwenden.
\begin{verbatim}
Kapitel(KID, Überschrift, Titel, Erscheinungsjahr, Name, Geburtsjahr)
\end{verbatim}
Beurteile diesen Vorschlag.

\item[]\hspace*{-\leftmargin}Gib jeweils eine SQL-Anfrage an, die die folgenden Ergebnisse liefert:

\item Name aller Autoren, die im Jahr 1970 geboren wurden.

\item Titel und Erscheinungsjahr aller Bücher, die vom Autor \enquote{J. K. Rowling}
veröffentlicht wurden, aufsteigend sortiert nach dem Erscheinungsjahr.

\item[]\hspace*{-\leftmargin}Das Datenbankschema soll so geändert werden, dass es auch möglich ist, einem Buch mehrere
Autoren zuzuordnen (z. B. bei Sammelbänden).

\item Erläutere, welche Änderungen an der Datenbank hierzu nötig sind.

\item Gib eine zu deinen Änderungen aus Teil f) passende SQL-Abfrage an, die die Titel
aller Bücher liefert, an denen der Autor \enquote{Sebastian Fitzek} beteiligt war.
\end{enumerate}

\begin{solutionbox}{19cm}
a) Primärschlüssel: Ein Attribut oder eine Attributkombination, das bzw. die jeden Datensatz
einer Relation eindeutig identifiziert.
Fremdschlüssel: Ein Attribut in einer Relation, das auf den Primärschlüssel einer anderen
Relation verweist, um Beziehungen zwischen Datensätzen herzustellen.

b) 

\includegraphics[width=8cm]{\pfad/cur11a.png}


c) Der Vorschlag ist nicht sinnvoll, da dadurch starke Redundanzen entstehen.
Bei jedem Kapitel würden Buchdaten und Autorendaten erneut gespeichert.
Das erhöht den Speicherbedarf und führt leichter zu Inkonsistenzen bei Änderungen.

d) 
\begin{verbatim}
SELECT Autor.Name 
  FROM Autor 
  WHERE Geburtsjahr = 1970
\end{verbatim}

e)
\begin{verbatim}
SELECT Buch.Titel, Erscheinungsjahr
    FROM Buch, Autor
    WHERE Buch.AID = Autor.AID
      AND Autor.Name = 'J. K. Rowling'
    ORDER BY Erscheinungsjahr ASC
\end{verbatim}

f)
Durch die Änderung entsteht eine n:m-Beziehung zwischen Buch und Autor.
Dazu wird der Fremdschlüssel AID aus der Tabelle Buch entfernt und eine
zusätzliche Beziehungstabelle eingeführt:

\texttt{Buch\_Autor(\fk{BID}, \fk{AID})}

g)
\begin{verbatim}
SELECT Buch.Titel
    FROM Buch, Buch_Autor, Autor
    WHERE Buch.BID = Buch_Autor.BID
      AND Buch_Autor.AID = Autor.AID
      AND Autor.Name = 'Sebastian Fitzek'
\end{verbatim}

\end{solutionbox}